% !TEX root = ../proyect.tex

\chapter{Introducción}\label{cap:introduccion}

\section{Marco conceptual}\label{sec:marco-conceptual}

% TODO: Describir el marco conceptual del proyecto
% - Contexto universitario actual
% - Problemática de acceso a información académica
% - Rol de los asistentes virtuales en educación

\section{Motivación}\label{sec:motivacion}

% TODO: Explicar la motivación del proyecto
% - Necesidad detectada en la Universidad de Sevilla
% - Beneficios esperados para estudiantes
% - Oportunidades tecnológicas (IA, NLP)

\section{Objetivos}\label{sec:objetivos}

\subsection{Objetivo general}\label{sec:objetivo-general}

% TODO: Objetivo principal del TFG

\subsection{Objetivos técnicos}\label{sec:objetivos-tecnicos}

% TODO: Objetivos técnicos específicos
% - Implementar asistente conversacional con Rasa
% - Integrar LLMs locales para Text-to-SQL
% - Diseñar base de datos académica escalable
% - etc.

\subsection{Objetivos académicos y personales}\label{sec:objetivos-academicos}

% TODO: Objetivos de aprendizaje
% - Profundizar en NLP y sistemas conversacionales
% - Aprender arquitecturas de IA
% - Desarrollar habilidades de gestión de proyectos
% - etc.

\section{Estructura de la memoria}\label{sec:estructura-memoria}

% TODO: Describir cómo está organizado el documento
